\section{What a Service Costs to Make}
\label{sec:ch12-what-a-service-costs-to-make}
\index{subgraph!what a new one costs|(}%

The domain code moved with \code{git mv} and changed in its namespace and using
lines and nowhere else, which is
chapter~\ref{ch:the-first-cut-extracting-catalog}'s finding repeated five times
and not worth a second section. What is worth a section is everything that had to be
written around it, because that is the real price of a service and nobody quotes
it.

Every one of the five needs a \code{csproj}, a \code{Program.cs}, a
\code{DbContext}, a data registration, a database seeder, a Dockerfile, a launch
profile, two \code{appsettings} files and a one-line \code{ModuleInfo.cs}, and
then a stanza in the compose file that is not its own. Ten files apiece that
contain no domain logic at all. Multiply by five and that is fifty files a
reader could be forgiven for thinking are the interesting part of a migration,
and they are the part that will take the longest and teach the least.

\index{Module attribute@\code{Module} attribute!one per assembly}%
The smallest of them is the one that explains the rest:

\begin{minted}[fontsize=\footnotesize]{csharp}
[assembly: Module("Pricing")]
\end{minted}

That attribute is assembly-scoped, which is why \code{Mosaic.Api} could collect
six domains behind one \code{AddMosaic()} and why there are six of these now,
each collecting one. It is the clearest single marker of what changed: the unit
of registration used to be the process and is now the domain.

\subsection*{Three subgraphs with no root field}

\index{Query type!a subgraph with no root field}%
Pricing, Inventory and Reviews contribute only fields to entities somebody else
owns. None of them has a root field, and that is the right shape rather than an
omission: the only way into those services is a representation the router
builds.

It is also the one thing in this chapter that failed at build time with an error
that names the wrong subject. All three refused to start:

\begin{minted}[fontsize=\footnotesize,breaklines]{text}
Unhandled exception: HotChocolate.SchemaException:

1. The schema builder was unable to identify the query type of the schema.
   Either specify which type is the query type or set the schema builder to
   non-strict validation mode.
\end{minted}

\code{AddApolloFederation()} puts \code{\_service} and \code{\_entities} on the
query type. It does not create one, and the source generator only emits one when
it finds a \code{[QueryType]} class. So a subgraph that is pure contribution
needs a line saying so:

\begin{minted}[fontsize=\footnotesize,breaklines]{csharp}
builder.AddGraphQL()
    .AddMosaicSubgraph()
    .AddQueryType()
    .AddPricing()
    .RegisterDbContextFactory<PricingDbContext>();
\end{minted}

Nothing in that error mentions federation, and the message a reader will search
for is about a query type rather than about a subgraph. Worth knowing before you
spend twenty minutes reading \code{AddApolloFederation} looking for the bug,
which is what I did.

\subsection*{The three settings that stopped being comments}

\index{subgraph defaults!in a shared library}%
Chapter~\ref{ch:the-first-cut-extracting-catalog} found three subgraph settings
that only fail once two schemas meet. They are the two cost options, and
\code{registerNodeInterface} set to false. At two services those were two copies
of a long comment.
At six they are a package, and the reason is a property of that class of setting
rather than of these three: a service that gets one of them wrong does not break
itself. It breaks everybody's composition, and the error names a type rather
than a service.

So they live in one method now, in the first project in this repository that is
not a service:

\begin{minted}[fontsize=\footnotesize,breaklines]{csharp}
public static IRequestExecutorBuilder AddMosaicSubgraph(this IRequestExecutorBuilder builder)
    => builder
        .AddApolloFederation()
        .AddGlobalObjectIdentification(registerNodeInterface: false)
        .ModifyCostOptions(options =>
        {
            options.ApplyCostDefaults = false;
            options.ApplySlicingArgumentDefaultValue = false;
        })
        .AddApplicationService<ILoggerFactory>()
        .AddDiagnosticEventListener<RequestTimelineListener>();
\end{minted}

The gate asserts that all six assemble the same thirteen middleware in the same
order. That check was chapter~\ref{ch:the-life-of-a-request}'s and it has a
second job now: six services calling one method should produce one pipeline, and
if one of them does not, its registrations have drifted from the platform's.

\subsection*{Catalog gets its instrumentation back}

\index{observability!six timelines are not a trace}%
Chapter~\ref{ch:the-first-cut-extracting-catalog} took the
\code{ServiceCallCounter} out of Catalog and gave a reason: the counter was the
monolith's instrumentation, and a service answering one domain's questions has
nothing to count with it. Chapter~\ref{ch:enter-the-router} then put the
consequence plainly: \enquote{Half of every federated query in this system is
now unmeasured, and the extraction that caused it did not warn me.}

That reason was true of two services and false of six. A graph in which one
process reports its costs and five do not is a graph where every question about
cost has the same answer, which is that it happened somewhere else. So all six
count now, and the counting is in the shared library so that the six answers
mean the same thing.

I want to be careful about how much that fixes, because it is less than it
looks. Six services each reporting their own timeline is six unrelated
timelines. Nothing in a response says which of them belong to the same federated
query, and joining them up is a different piece of work rather than a bigger
version of this one. Chapter~\ref{ch:observability} owns it and this chapter has
not done it.

\medskip
That is the boilerplate. The genuinely difficult decision in a split this size
is not which files to write but which ones to write twice.

\index{subgraph!what a new one costs|)}%
